\documentclass[11pt,a4paper]{article}
\usepackage{amsmath,amssymb,bm}
\usepackage[margin=2.5cm]{geometry}
\usepackage{booktabs}
\usepackage{hyperref}

\title{The chiral colour-dielectric model \\
\large as implemented in \texttt{eos.ccdm}}
\author{}
\date{}

\newcommand{\phib}{\bar\varphi}
\newcommand{\Mstar}{M^{*}}
\newcommand{\mustar}{\mu^{*}}

\begin{document}
\maketitle

\begin{abstract}
A mean-field model of deconfined $u,d,s$ matter in which confinement and
chiral symmetry breaking are two faces of one mechanism. A dilaton field
carries the gluon condensate; the dielectric function built from it measures
the medium's transparency to colour and divides the quark masses, so that as
the condensate reaches its vacuum value the effective masses diverge and the
quarks leave the medium entirely. This document states every equation the code
solves and every quantity it returns. The implementation specification is
\texttt{docs/ccdm\_implementation.md} and is the authority wherever it and this
document differ, with the two exceptions recorded in \S\ref{sec:corrections},
which are places where the specification contradicts the thermodynamic audit
it itself mandates.
\end{abstract}

\section{Conventions}

Natural units $\hbar=c=k_B=1$ throughout the physics; masses, chemical
potentials and the dilaton in MeV, densities in MeV$^3$, $\Omega$, $P$ and
$\varepsilon$ in MeV$^4$. The public boundary of the code is fm-based, with
$\hbar c = 197.3269804$ MeV\,fm.

Flavour $f\in\{u,d,s\}$, colour $a\in\{r,g,b\}$, and the nine colour--flavour
modes $j=(f,a)$ indexed flavour-major. Electric charges $q_u=+\tfrac23$,
$q_d=q_s=-\tfrac13$. \textbf{Strangeness is $S=+1$ per $s$ quark}, the opposite
of the PDG sign and the convention used consistently in this repository.
$C$ is the electric charge of strongly interacting matter only: the leptons are
excluded from it and enter through the separate condition of total electric
neutrality.

Colour generators are taken in the normalisation
\begin{equation}
(T_3)_{r,g,b}=\left(+\tfrac12,-\tfrac12,0\right),
\qquad
(T_8)_{r,g,b}=\left(+\tfrac13,+\tfrac13,-\tfrac23\right),
\end{equation}
i.e. $T_8=\lambda_8/\sqrt3$. Three normalisations are in circulation in the
colour-superconductivity literature and mixing them corrupts $\mu_8$ by factors
between $1.15$ and $1.7$; the conversions are documented in
\texttt{eos.general.pairing}.

\section{The Lagrangian}

Four boson fields --- the dilaton $\varphi$, the light scalar $\sigma$, the
pion triplet $\bm\pi$ and the strange scalar $\zeta$ --- plus the vector meson
$\omega_\mu$ and the diquark interaction:
\begin{align}
\mathcal L &= \sum_f \bar q_f\left[i\gamma^\mu\partial_\mu - \Mstar_f\right]q_f
 + \tfrac12(\partial\varphi)^2 - U(\varphi)
 + \tfrac12(\partial\sigma)^2 + \tfrac12(\partial\bm\pi)^2
 + \tfrac12(\partial\zeta)^2 - V(\sigma,\bm\pi,\zeta) \nonumber\\
&\quad + \mathcal L_{\rm vec} + \mathcal L_{\rm pair},
\\
\mathcal L_{\rm vec} &= -g_\omega\,\bar q\gamma^\mu q\,\omega_\mu
 - \tfrac14 F_{\mu\nu}F^{\mu\nu} + \tfrac12 m_\omega^2\omega_\mu\omega^\mu,
\\
\mathcal L_{\rm pair} &= G_D\sum_{\eta=1}^{3}
 \left|\bar q\,i\gamma_5 C\,\epsilon^{ab\eta}\epsilon_{ij\eta}\,\bar q^{\,T}\right|^2,
\qquad \eta=1,2,3 \leftrightarrow (ds),(us),(ud).
\end{align}

\subsection{The dielectric, and why the fourth power}

The dielectric function is \emph{not} a field: it has no kinetic term, no
potential and no field equation. It is an abbreviation,
\begin{equation}
\boxed{\;\chi(\varphi) = \left[1-\phib^{\,4}\right]^{p},
\qquad \phib \equiv \varphi/\varphi_0,\qquad p=1 \;}
\label{eq:chi}
\end{equation}
and it appears only in the quark mass operator. The fourth power is fixed:
$\phib^{\,4}$, not $\phib$, is the gluon condensate, because the dilaton is a
canonically normalised scalar of mass dimension one while
$\langle G^A_{\mu\nu}G^{A\mu\nu}\rangle$ has dimension four. The scale anomaly
makes this sharp --- the combination the anomaly fixes comes out a pure power
of $\phib$ only for exponent four,
\begin{equation}
4U - \varphi\,\frac{\partial U}{\partial\varphi} = 4B_g\left(1-\phib^{\,4}\right),
\end{equation}
whereas $\phib^{2}$, $\phib^{3}$, $\phib^{5}$ or $\phib^{6}$ leave a residual
$\phib^{\,n}\ln\phib$. So
\begin{equation}
1-\phib^{\,4} = 1 - \frac{\langle G^2\rangle_{\rm med}}{\langle G^2\rangle_{\rm vac}},
\end{equation}
and the medium becomes transparent to colour exactly in proportion to how much
condensate has melted. The two endpoints are then correct by construction:
$\chi\to0$ at $\phib=1$ (confinement, $\Mstar\to\infty$) and $\chi\to1$ at
$\phib=0$ (perturbative, $\Mstar\to m_f$). The exponent $p$ outside is a
separate matter: only $\chi^{p}$ enters $\Mstar$, so $p$ and the bracket are
meaningful only as a pair, and the code locks $p=1$.

\subsection{The solve variable}
\label{sec:Phi}

Define
\begin{equation}
\Phi \equiv \phib^{\,4}
= \frac{\langle G^2\rangle_{\rm med}}{\langle G^2\rangle_{\rm vac}},
\end{equation}
in which the potential and the dielectric are simply
\begin{equation}
U = B_g\left[\Phi(\ln\Phi-1)+1\right],
\qquad
\chi = (1-\Phi)^{p},
\qquad
\frac{\mathrm dU}{\mathrm d\Phi} = B_g\ln\Phi .
\label{eq:U}
\end{equation}

\textbf{$\varphi$ is the field; $\Phi$ is the solve variable.} $\Phi$ cannot be
the Lagrangian field: its kinetic term
$\tfrac12(\partial\varphi)^2 = \tfrac{\varphi_0^2}{32}\Phi^{-3/2}(\partial\Phi)^2$
is not canonical and diverges at $\Phi\to0$, $[\Phi]=4$ makes
$(\partial\Phi)^2$ dimension ten so no renormalisable Lagrangian exists in it,
and the glueball mass would stop being $U''$. But as a solve variable it is
strictly better, for one concrete reason: written in $\phib$ the dilaton
residual has a \emph{spurious root at $\phib=0$}, where both of its terms vanish
as $\phib^{\,3}$ --- the first as $\phib^{\,3}\ln\phib$, the second as
$\phib^{\,3}$ --- so it vanishes there for \emph{any} scalar density. It is an
artefact of the parametrisation: the Jacobian $\mathrm d\Phi/\mathrm d\phib =
4\phib^{\,3}$ vanishes, not the physics. In $\Phi$, where
$\mathrm dU/\mathrm d\Phi=B_g\ln\Phi\to-\infty$, no such root exists.

\subsection{Mean fields}

At homogeneous mean field every boson is replaced by its expectation value,
$\bm\pi\to0$ by parity, and $\omega_\mu\to\delta_{\mu0}\omega_0$:
\begin{equation}
\mathcal L_{\rm MF} = \sum_f \bar q_f\left[i\gamma^\mu\partial_\mu
 - g_\omega(n_B)\gamma^0\omega_0 - \Mstar_f\right]q_f
 + \mathcal L^{\rm MF}_{\rm pair}
 - \left[U(\varphi)+V(\sigma,\zeta) - \tfrac12 m_\omega^2\omega_0^2\right],
\end{equation}
with
\begin{equation}
\boxed{\;
\Mstar_u = \frac{g_q\sigma+m_u}{\chi},\qquad
\Mstar_d = \frac{g_q\sigma+m_d}{\chi},\qquad
\Mstar_s = \frac{g_s\zeta+m_s}{\chi}\;}
\label{eq:masses}
\end{equation}
and the chiral potential
\begin{align}
V(\sigma,\zeta) &= \frac{\lambda}{4}\left(\sigma^2-v^2\right)^2
 + \frac{\lambda_\zeta}{4}\left(\zeta^2-v_\zeta^2\right)^2
 - \epsilon_\sigma\sigma - \epsilon_\zeta\zeta + C_0,
\label{eq:V}\\
\frac{\partial V}{\partial\sigma} &= \lambda\sigma(\sigma^2-v^2)-\epsilon_\sigma,
\qquad
\frac{\partial V}{\partial\zeta} = \lambda_\zeta\zeta(\zeta^2-v_\zeta^2)-\epsilon_\zeta .
\label{eq:dV}
\end{align}
$v$ and $v_\zeta$ are constants, not fields: the Mexican-hat radii before
explicit breaking, and not observables. The linear terms shift the true minima
to $\sigma_0,\zeta_0$.

\subsection{The diquark condensate}

The three condensates
$s_\eta = \langle q^{T}C\gamma_5\,\epsilon^{ab\eta}\epsilon_{ij\eta}\,q\rangle$
give, to quadratic order in the fluctuation,
\begin{equation}
\mathcal L^{\rm MF}_{\rm pair} = \sum_{\eta=1}^{3}
 \left[\frac{\Delta_\eta^*\mathcal O_\eta + \Delta_\eta\mathcal O_\eta^\dagger}{2}
 - \frac{|\Delta_\eta|^2}{4G_D}\right],
\qquad
\Delta_\eta = 2G_D s_\eta .
\end{equation}
The mean-field treatment is exact for an auxiliary field appearing
quadratically, so substituting $\Delta_\eta=2G_Ds_\eta$ back reproduces
$G_D|s_\eta|^2$ exactly; the gap equation
$\partial\Omega/\partial\Delta_\eta=0$ \emph{is} the condensate definition; and
at the solution $\Omega_\eta=-G_D|s_\eta|^2<0$, so condensation lowers the
grand potential as it must. All $\Delta_\eta$ are taken real, which is a
convention since the phases are unobservable in a homogeneous condensate.

\textbf{The sign of each $\Delta_\eta$ is a gauge.} $\Omega$ is invariant under
flipping any subset of the three gaps and each gap kernel flips with its own
gap, so $-\Delta_\eta$ is a root whenever $\Delta_\eta$ is. The code reports
magnitudes.

\section{Chemical potentials}

In the conserved-charge basis, with the colour potentials present only where
pairing is,
\begin{equation}
\mu_{f,a} = \tfrac13\mu_B + q_f\mu_C + s_f\mu_S + (T_3)_a\mu_3 + (T_8)_a\mu_8 .
\label{eq:mu}
\end{equation}
Both colour generators are traceless, so colour drops out of $n_B$. The
potential fed to every momentum integral is shifted by the vector self-energy,
\begin{equation}
\boxed{\;\mustar_{f,a} = \mu_{f,a} - \Sigma_V,
\qquad \Sigma_V = g_\omega(n_B)\,\omega_0 + \Sigma_R\;}
\label{eq:mustar}
\end{equation}
with $\Sigma_R$ the rearrangement term of \S\ref{sec:vector}.

$\mu_C$ is the \emph{charge} chemical potential and $\mu_e=-\mu_C$, so beta
equilibrium reads $\mu_C+\mu_e=0$. Writing $+q_f\mu_e$ in place of $+q_f\mu_C$
flips the $u$--$d$ ordering: at $\mu_B=1200$, $\mu_C=-53.15$ MeV the correct
values are $\mu_u=364.6<\mu_d=417.7$ MeV, whereas the wrong sign gives
$\mu_u=435.4>\mu_d=382.3$ MeV.

Leptons are free Fermi gases at $\mu_\ell=\mu_{L_\ell}-\mu_C$, with masses
$m_e=0.510999$ and $m_\mu=105.658$ MeV and degeneracy $g=2$. Trapped neutrinos
carry $g=1$, not $2$: they are left-handed only.

\section{The vector coupling and its rearrangement}
\label{sec:vector}

The vector coupling is a function of the state, not a parameter:
\begin{equation}
g_\omega(n_B) = \frac{\bar g_\omega}{1+(n_B/n_c)^2},
\qquad
\frac{\partial g_\omega}{\partial n_B}
 = -\frac{2\bar g_\omega\,n_B/n_c^2}{\left[1+(n_B/n_c)^2\right]^2}.
\label{eq:gomega}
\end{equation}
Its field equation is sourced by the \textbf{quark} number density
$n_q=3n_B$, because the coupling in $\mathcal L_{\rm vec}$ is to
$\bar q\gamma^\mu q$; using $n_B$ understates $\omega_0$ by a factor three and
the repulsive energy by nine:
\begin{equation}
\omega_0 = \frac{g_\omega(n_B)\,n_q}{m_\omega^2}.
\label{eq:omega}
\end{equation}
Because $g_\omega$ depends on the density, the shift of $\mu$ is the derivative
of the interaction energy $W=\tfrac12 m_\omega^2\omega_0^2$ with respect to the
quark density, not merely $g_\omega\omega_0$:
\begin{equation}
\boxed{\;
\Sigma_V = \frac{\mathrm dW}{\mathrm dn_q}
 = g_\omega\omega_0 + \Sigma_R,
\qquad
\Sigma_R = \frac{\partial g_\omega}{\partial n_B}\,\omega_0\,n_B \;}
\label{eq:SigmaR}
\end{equation}
(the chain rule through $n_B=n_q/3$ is what turns the naive
$(\partial g/\partial n_q)n_q$ into the $n_B$ written here).
\textbf{$\Sigma_R$ enters $\mu$ and $P$ and never $\varepsilon$.} Note that
$W$ peaks at $n_B=n_c$ and falls beyond it, so $\Sigma_R$ --- and with it
$\Sigma_V$ --- changes sign there; this is a property of the coupling form
\eqref{eq:gomega} and not of the implementation.

The diquark coupling may be dressed by the dielectric,
\begin{equation}
G_D \longrightarrow G_D/\chi^{\,q},\qquad q\in\{0,1\},
\end{equation}
$q=1$ being the exponent a gluon-exchange origin gives and the largest that
leaves the pairing channel confined, the criterion being $q\le p$ because the
critical coupling for vacuum diquark condensation grows only linearly in
$\Mstar\propto\chi^{-p}$.

\section{The ideal-gas integrals}

Every momentum integral in this model is a relativistic ideal gas of one
colour--flavour mode with spin degeneracy $g=2$. With
$E_k=\sqrt{k^2+{\Mstar}^2}$ and
\begin{equation}
f^{\pm}_k = \left[1+e^{(E_k\mp\mustar)/T}\right]^{-1},
\end{equation}
the five integrals are
\begin{align}
\mathcal N(\Mstar,\mustar,T) &= \frac{g}{2\pi^2}\int_0^{k_{\max}}\!\!\mathrm dk\;k^2
 \left[f^+_k-f^-_k\right],
&&\text{(antiparticles \emph{subtract})}
\label{eq:N}\\
\mathcal R_s(\Mstar,\mustar,T) &= \frac{g}{2\pi^2}\int_0^{k_{\max}}\!\!\mathrm dk\;k^2\,
 \frac{\Mstar}{E_k}\left[f^+_k+f^-_k\right],
&&\text{(antiparticles \emph{add})}
\label{eq:Rs}\\
\mathcal P(\Mstar,\mustar,T) &= \frac{g}{2\pi^2}\int_0^{k_{\max}}\!\!\mathrm dk\;k^2\,T
 \sum_{\pm}\ln\left(1+e^{-(E_k\mp\mustar)/T}\right),
\label{eq:P}\\
\mathcal E(\Mstar,\mustar,T) &= \frac{g}{2\pi^2}\int_0^{k_{\max}}\!\!\mathrm dk\;k^2 E_k
 \left[f^+_k+f^-_k\right],
\label{eq:E}\\
\mathcal S(\Mstar,\mustar,T) &= \frac{g}{2\pi^2}\int_0^{k_{\max}}\!\!\mathrm dk\;k^2
 \sum_{\pm}\left[\frac{E_k\mp\mustar}{T}f^\pm_k
 + \ln\left(1+e^{-(E_k\mp\mustar)/T}\right)\right].
\label{eq:S}
\end{align}
They satisfy, per mode and to machine precision,
\begin{equation}
\mathcal N = \frac{\partial\mathcal P}{\partial\mustar},\quad
\mathcal R_s = -\frac{\partial\mathcal P}{\partial\Mstar},\quad
\mathcal S = \frac{\partial\mathcal P}{\partial T},\quad
\mathcal E = -\mathcal P + \mustar\mathcal N + T\mathcal S .
\label{eq:modeeuler}
\end{equation}
Note that the single-mode Euler relation carries $T\mathcal S$ and \emph{not}
$\Mstar\mathcal R_s$: the scalar density is a response, not a conserved charge.

\textbf{The medium integrals are unregularised.} $k_{\max}$ is a numerical
ceiling at which the integrand has died,
$k_{\max}=\max(|\mustar|,\Mstar)+45T+12\Mstar+200$ MeV, and not a parameter;
the cutoff $\Lambda$ of \S\ref{sec:pairing} applies to the \emph{pairing}
integral alone. This is why a sharp cutoff is admissible here, and it implies
the validity ceiling $\mu_{\max}=\sqrt{\Lambda^2+\Mstar{}^2}$ for the pairing
sector, which the code declares rather than exceeds.

\subsection{Quadrature, and confinement}

The integrals are evaluated by \textbf{panel-split Gauss--Legendre} quadrature
with breakpoints at the Fermi momentum $k_F=\sqrt{{\mustar}^2-{\Mstar}^2}$ and at
$k_F\pm 25T$. This is not a refinement: a single panel cannot resolve a step of
width $\sim T$ inside a range of $\sim1500$ MeV, and single-panel accuracy is
not even monotone in the node count, because whether a node lands near the
Fermi step is accidental. At $T=0.5$ MeV the split scheme reaches with $100$
nodes per panel what a single panel needs $5000$ to reach.

\textbf{At $T=0$ a mode with $\mustar\le\Mstar$ contributes identically zero.}
All five integrals \eqref{eq:N}--\eqref{eq:S} return exactly $0$, not a small
number. This is the confinement mechanism itself: as $\phib\to1$ the dielectric
closes, $\Mstar$ diverges through \eqref{eq:masses}, and the quarks leave the
medium. Smoothing the threshold destroys the pinning that makes the
deconfinement transition first order.

At $T=0$ and $\mustar>\Mstar$ the closed forms are
\begin{align}
\mathcal N &= \frac{g}{6\pi^2}k_F^3,
\qquad
\mathcal R_s = \frac{g}{4\pi^2}\Mstar\left(k_F\mustar
 - \Mstar{}^2\ln\frac{k_F+\mustar}{\Mstar}\right),
\qquad \mathcal S = 0,
\\
\mathcal P &= \frac{g}{48\pi^2}\left[2k_F^3\mustar
 - 3\Mstar{}^2\left(\mustar k_F - \Mstar{}^2\ln\frac{k_F+\mustar}{\Mstar}\right)\right],
\\
\mathcal E &= \frac{g}{16\pi^2}\left[2k_F^3\mustar
 + \Mstar{}^2\left(\mustar k_F - \Mstar{}^2\ln\frac{k_F+\mustar}{\Mstar}\right)\right].
\end{align}

\section{The pairing sector}
\label{sec:pairing}

The pairing machinery is shared with the NJL model of this repository and is
stated here in full so that this document is self-contained.

\subsection{The gap matrix and the BdG problem}

Build the $9\times9$ gap matrix on the colour--flavour modes directly from the
antisymmetric tensors,
\begin{equation}
\mathcal G = \sum_{\eta=1}^{3}\Delta_\eta\,\mathcal B_\eta,
\qquad
\left(\mathcal B_\eta\right)_{(ai),(bj)} = \epsilon^{ab\eta}\epsilon_{ij\eta}.
\label{eq:gapmatrix}
\end{equation}
Its eigenvalue multiplicities are \emph{derived}, never assigned: unpaired
$0^{(\times9)}$; 2SC $\pm\Delta^{(\times2)},0^{(\times5)}$; uSC/dSC
$\pm\Delta^{(\times2)},\pm\sqrt{\Delta_a^2+\Delta_b^2},0^{(\times3)}$; CFL
$\Delta^{(\times3)},-\Delta^{(\times5)},2\Delta$.

\textbf{The eigenvalues alone are not enough.} As soon as the quark masses
differ,
$\left[\mathcal G,\ \mathrm{diag}(\Mstar_f)\right]\neq0$, so the gap matrix and
the mass matrix share no eigenbasis and the problem does not factorise into
independent scalar dispersions. At $\Mstar=(60,65,300)$, $\mustar=450$,
$\Delta=60$ MeV the eigenvalue prescription misses branches by up to $33.6$ MeV
in uSC and $27.1$ in CFL. So at each momentum the code assembles and
diagonalises
\begin{equation}
\boxed{\;
H_{\rm BdG}(k) = \begin{pmatrix} \xi(k) & \mathcal G \\ \mathcal G & -\xi(k)\end{pmatrix},
\qquad
\xi(k) = \mathrm{diag}\left(E_{k,f(j)} - r\,\mustar_j\right)_{j=1}^{9},\;}
\label{eq:bdg}
\end{equation}
an $18\times18$ real symmetric matrix, once for particles ($r=+1$) and once for
antiparticles ($r=-1$). The nine quasiparticle energies $E^{(j),r}_\Delta(k)$
are the \textbf{non-negative half of the signed spectrum},
$\mathrm{sort}(\mathrm{eigvalsh})[9{:}]$, and not the nine largest in modulus:
the two agree in value, but only the first is smooth through a gapless window,
and that smoothness is what makes the Hellmann--Feynman derivatives below carry
the correct branch sign.

\subsection{The pairing potential, as a correction}

\begin{align}
\Omega_{\rm quark} &= -\sum_{j=1}^{9}\mathcal P\!\left(\Mstar_j,\mustar_j,T\right)
 + \delta\Omega_{\rm pair} + \sum_{\eta=1}^{3}\frac{\Delta_\eta^2}{4G_D},
\label{eq:Omegaquark}\\
\delta\Omega_{\rm pair} &= -\frac{1}{2\pi^2}\int_0^{\Lambda}\!\!\mathrm dk\,k^2
 \sum_{r=\pm}\sum_{j=1}^{9}
 \Bigl\{\left(E^{(j),r}_\Delta - \left|\xi^{(j),r}\right|\right)
 + 2T\Bigl[\ln\bigl(1+e^{-E^{(j),r}_\Delta/T}\bigr)
 - \ln\bigl(1+e^{-|\xi^{(j),r}|/T}\bigr)\Bigr]\Bigr\}.
\label{eq:dOmega}
\end{align}
Written as a \emph{correction} --- a difference from the unpaired spectrum ---
every brace vanishes identically at $\Delta_\eta=0$, so $\delta\Omega_{\rm
pair}=0$ there exactly and the whole expression reduces to the unpaired one
mode by mode. The alternative "replacement" form, in which the paired modes'
Fermi integrals are dropped and the $|E|$ integral put in their place, differs
by a $\Delta$-independent but $\Lambda$-dependent constant of about 2\% of the
term: the gap comes out right and the equation of state is wrong.

\textbf{The antiparticle branches are not optional.} At $T=0$ they contribute
$8.8\%$ of the pairing potential at $\Lambda=600$ MeV and $17.1\%$ at
$\Lambda=1000$ MeV, so omitting them is a cutoff-dependent error that cancels
nowhere. The thermal logarithm is computed as
$2T\,\mathrm{logaddexp}(0,-|E|/T)$, never as $2T\ln(1+e^{-|E|/T})$, which
overflows for $|E|/T\gtrsim700$.

\subsection{The three gap equations}

There is one residual per channel, not one gap:
\begin{equation}
\boxed{\;
R_{\Delta_\eta} = \frac{\Delta_\eta}{2G_D}
 - \frac{1}{2\pi^2}\int_0^{\Lambda}\!\!\mathrm dk\,k^2\sum_{r=\pm}\sum_{j=1}^{9}
 \left\langle \mathcal V^{(j),r}\left|
 \begin{pmatrix} 0 & \mathcal B_\eta \\ \mathcal B_\eta & 0\end{pmatrix}
 \right|\mathcal V^{(j),r}\right\rangle
 \tanh\frac{E^{(j),r}_\Delta}{2T} = 0\;}
\label{eq:gapeq}
\end{equation}
with $\mathcal V^{(j),r}$ the eigenvector of $H_{\rm BdG}$ belonging to
$E^{(j),r}_\Delta$. The matrix element is the Hellmann--Feynman derivative
$\partial E^{(j),r}_\Delta/\partial\Delta_\eta$, and
$\partial H_{\rm BdG}/\partial\Delta_\eta$ is the constant matrix
$\mathcal B_\eta$ in the off-diagonal blocks, assembled once.

\textbf{This is never the $\Delta/|E|$ form}, which is obtained by
differentiating $|E|$ as though every branch were positive. That form is wrong
by a factor $12$ in the gapless window, and it makes the gap \emph{grow} with
the Fermi-surface mismatch --- the opposite of the physics.

\subsection{Paired densities and entropy}

For the modes the gap touches, the densities are \emph{not} the unpaired Fermi
integrals: pairing redistributes occupation around the Fermi surface, so
\begin{equation}
n_j = -\frac{\partial\Omega_{\rm quark}}{\partial\mu_j},
\qquad
\rho_{s,j} = -\frac{\partial\Omega_{\rm quark}}{\partial\Mstar_j},
\end{equation}
taken by Hellmann--Feynman on \eqref{eq:bdg} in the same quadrature pass.
Substituting the unpaired formula is a $15\%$ density error at a realistic gap
and breaks neutrality and Euler simultaneously. The entropy is built from the
quasiparticle occupations $n^{(j),r}_k=[1+e^{E^{(j),r}_\Delta/T}]^{-1}$; using
the unpaired entropy in a gapped phase is a four-orders-of-magnitude error at
low $T$. The suppression is a property of the \emph{spectrum}: at matched Fermi
surfaces the entropy ratio is $2.0\times10^{-4}$ at $T=5$ MeV and
$\Delta=60$ MeV, while at $\Mstar_s=300$ MeV the mismatch pushes the lowest
branch from $60$ to $11.9$ MeV and the ratio only to $0.20$.

Colour neutrality is imposed \emph{within} the pattern:
\begin{equation}
n_3 = \sum_f\left(n_{(f,r)}-n_{(f,g)}\right)=0,
\qquad
n_8 = \sum_f\left(n_{(f,r)}+n_{(f,g)}-2n_{(f,b)}\right)=0 .
\label{eq:colour}
\end{equation}
In an \emph{unpaired} region $n_3$ and $n_8$ vanish identically at
$\mu_3=\mu_8=0$ whatever else the state does, so the colour potentials are
unconstrained there and are pinned rather than solved for; a paired phase's
$\mu_8$ cannot be inherited from an unpaired solution.

\section{Assembly}
\label{sec:assembly}

Write the mode sums
$S_{\mathcal P}=\sum_j\mathcal P_j$, $S_{\mathcal E}=\sum_j\mathcal E_j$,
$S_{\mathcal S}=\sum_j\mathcal S_j$, the pairing cost
$D=\sum_\eta\Delta_\eta^2/4G_D$, and the vector field energy
$W=\tfrac12 m_\omega^2\omega_0^2$. Then
\begin{align}
\boxed{\;\Omega = U(\Phi) + V(\sigma,\zeta) - W - \Sigma_R\,n_q
 - S_{\mathcal P} + \delta\Omega_{\rm pair} + D
 - \sum_\ell\mathcal P_\ell - \mathcal P_\gamma\;}
\label{eq:Omega}\\
\boxed{\;\varepsilon = S_{\mathcal E} + U + V + W + \varepsilon_{\rm pair}
 + \sum_\ell\mathcal E_\ell + \mathcal E_\gamma\;}
\label{eq:eps}
\end{align}
with $P=-\Omega$,
\begin{equation}
\varepsilon_{\rm pair} = \delta\Omega_{\rm pair} + D
 + T\,\delta s_{\rm pair} + \sum_j \mustar_j\,\delta n_j,
\qquad
s = S_{\mathcal S} + \delta s_{\rm pair}
 + \sum_\ell\mathcal S_\ell + \tfrac{4}{45}\pi^2T^3,
\end{equation}
$\mathcal P_\gamma=\pi^2T^4/45$ and $\mathcal E_\gamma=3\mathcal P_\gamma$.
The conserved-charge sums are
\begin{equation}
n_f = \sum_a n_{(f,a)},\qquad
n_B = \tfrac13\sum_f n_f,\qquad
n_C = \sum_f q_f n_f,\qquad
n_S = n_s .
\end{equation}

\textbf{The signs.} $U+V$ enter $\varepsilon$ positively and $P$ negatively;
the vector field energy $W$ enters \emph{both} positively; the rearrangement
term $\Sigma_R n_q$ enters $P$ only. No vacuum subtraction appears anywhere,
because $U$ and $V$ both vanish at the physical vacuum by construction
(\S\ref{sec:parameters}).

The audit that must hold at every solved point, and does to machine precision:
\begin{equation}
\boxed{\;\varepsilon + P = Ts + \sum_j \mu_j n_j\;}
\label{eq:euler}
\end{equation}
together with $f=\varepsilon-Ts=-P+\sum_j\mu_jn_j$ and $P=-\Omega$ computed
from the two assemblies independently.

\subsection{Two corrections to the specification}
\label{sec:corrections}

\texttt{docs/ccdm\_implementation.md} is the authority for this model, and two
places in its assembly do not survive the audit \eqref{eq:euler} that its own
\S9.6 mandates.

\begin{enumerate}
\item Its \S4.3 writes $\varepsilon$ with $-\tfrac12 m_\omega^2\omega_0^2$.
\textbf{The sign must be $+$.} A repulsive vector interaction adds energy
density: the Hamiltonian density of the mean vector field is
$g_\omega\omega_0 n_q - \tfrac12 m_\omega^2\omega_0^2 = +\tfrac12
m_\omega^2\omega_0^2$ once \eqref{eq:omega} is used. This is the standard
density-dependent mean-field result, in which the scalar potentials enter
$\varepsilon$ positively and $P$ negatively while the vector term enters both
positively.

\item Its \S4.1 carries $\Sigma_R$ in $\mustar$ but omits the compensating
$-\Sigma_R n_q$ from $\Omega$. Without it, $n=-\partial\Omega/\partial\mu$ and
\eqref{eq:euler} both fail as soon as $g_\omega$ depends on the density. The
rearrangement term belongs in $\mu$ and $P$ and never in $\varepsilon$.
\end{enumerate}

With either error present the Euler residual is of order percent while every
other quantity still looks like a reasonable equation of state; with both
corrections it is $\sim10^{-16}$.

\section{The residual}
\label{sec:residual}

The unknown vector, in the order the solver assembles it, is
\begin{equation}
\vec X = \bigl(\Phi,\ \sigma,\ \zeta,\ [\Sigma_V],\
 [\Delta_\eta]_{\eta\in\rm free},\ [\mu_3,\mu_8],\
 \mu_B,\ \mu_C,\ [\mu_S],\ [\mu_{\nu_e}]\bigr),
\end{equation}
the bracketed entries being present only where the model or the mode has them:
$\Sigma_V$ where $\bar g_\omega\neq0$, the gaps as the pairing pattern
declares, the colour potentials whenever the pattern pairs at all, $\mu_S$ iff
the mode holds $Y_S$, and $\mu_{\nu_e}$ iff it holds $Y_{L_e}$.

The rows, in the same order:
\begin{align}
R_1 &= B_g\ln\Phi + \frac{p}{1-\Phi}\sum_f \Mstar_f\,\rho_{s,f} &&\text{dilaton}
\label{eq:R1}\\
R_2 &= \frac{\partial V}{\partial\sigma}
 + \frac{g_q}{\chi}\left(\rho_{s,u}+\rho_{s,d}\right) &&\text{light scalar}
\label{eq:R2}\\
R_3 &= \frac{\partial V}{\partial\zeta}
 + \frac{g_s}{\chi}\,\rho_{s,s} &&\text{strange scalar}
\label{eq:R3}\\
R_4 &= \Sigma_V - \left[g_\omega(n_B)\,\omega_0 + \Sigma_R\right] &&\text{vector}
\label{eq:R4}\\
R_{\Delta_\eta} &= \text{Eq.~\eqref{eq:gapeq}} &&\text{one per free gap}\\
R_{n_3},\ R_{n_8} &= \text{Eq.~\eqref{eq:colour}} &&\text{colour neutrality}\\
R_{n_B} &= n_B - n_B^{\rm target} &&\text{the density}\\
&\quad\text{+ the mode's charge rows (\S\ref{sec:modes})}. \nonumber
\end{align}

Equation \eqref{eq:R1} is $\partial\Omega/\partial\Phi$, obtained from
$\partial\Omega/\partial\phib$ by dividing out the Jacobian $4\phib^{\,3}$;
that division is exactly what removes the spurious root of
\S\ref{sec:Phi}. Equations \eqref{eq:R1}--\eqref{eq:R3} are minima of $\Omega$,
with the boundary minimiser $\sigma\to0$ admitted --- that is how chiral
restoration appears.

\textbf{$R_4$ carries $\Sigma_V$, not $\omega_0$.} The vector field is
circular: $\omega_0$ sets $\mustar$, which sets $n_B$, which sets $\omega_0$.
Carried as the total shift $\Sigma_V$ everything downstream is explicit ---
$\mustar$ from $\Sigma_V$, the densities from $\mustar$, $\omega_0$ and
$\Sigma_R$ from the densities --- and \eqref{eq:R4} is the single statement
that the returned field is the one the returned densities source. This is also
this repository's declared convention for a density-dependent coupling: the
unknown vector uses the effective potentials, so the rearrangement and the
large vector shift cancel out of the iteration.

Each row is divided by the scale of the quantity it balances before the
residual norm is taken: $R_1$ by $B_g$, $R_2$ and $R_3$ by $\epsilon_\sigma$
and $\epsilon_\zeta$, $R_4$ by $\mu_B/3$, the gap and colour rows by
$(\mu_B/3)^3/\pi^2$, and the density and charge rows by $n_B$. Unscaled, the
rows span twenty-two orders of magnitude and the norm would report whichever
carries the largest units.

\section{Branches and patterns: what is enumerated}

Neither the chiral/dielectric \textbf{branch} nor the pairing \textbf{pattern}
is a mode: both are decided by which candidate minimises the free energy.

Below the deconfinement onset two chiral branches coexist at fixed dilaton ---
a confined one, where $\sigma\simeq f_\pi$, the dielectric is nearly opaque and
the quarks are too heavy to appear at all, and a restored one, where $\sigma$
has collapsed and the quarks are present. \textbf{A solver that alternates
between updating $\sigma$ and $\omega_0$ two-cycles between them} and exits
with a mixed state --- $\sigma$ from one branch, $\omega_0$ from the other ---
which reads as a spuriously deep minimum at zero quark density. So each branch
is seeded separately, solved to self-consistency, and compared by $\Omega$; a
branch that fails to converge is reported missing, never replaced by a
neighbour, because substituting a converged neighbouring point is how a fake
phase boundary appears.

\begin{table}[h]
\centering
\begin{tabular}{lll}
\toprule
branch & seed $(\Phi,\sigma,\zeta)$ & \\
\midrule
confined & $(1^-,\ f_\pi,\ \zeta_0)$ & the vacuum: $n_B=0$, $P=0$ exactly \\
restored & $(0.4^4,\ 0,\ 0)$ & \\
partial & $(0.4^4,\ 0.5f_\pi,\ 0.7\zeta_0)$ & \\
\bottomrule
\end{tabular}
\end{table}

The confined branch is enumerated at fixed \emph{potential}, where it is what
the deconfined branch must beat and the crossing is the deconfinement onset,
and \emph{not} at fixed density: with the dielectric closed $n_B=0$
identically, so no nonzero density row can be met and the Jacobian's field
columns all vanish.

The pairing patterns --- unpaired, 2SC, uSC, dSC, CFL, and one asymmetric free
seed --- are declarations of which $\Delta_\eta$ are unknowns. The gap equation
has three roots at any Fermi-surface mismatch (zero, a barrier maximum, and the
physical BCS root), so which root a solve lands on is decided by the seed;
enumerating the seeds is what makes the answer an answer rather than one solve
repeated. Gapless states, $\min_k E^{(j)}_\Delta<0$, are \emph{flagged} rather
than ranked: they are not stationary points of the same functional and
comparing $\Omega$ across one is not valid.

The candidate set of a fixed-density solve is the \emph{product} of the two
enumerations, because which pattern survives depends on the strange quark's
effective mass, which is a property of the branch.

\textbf{The first-order transition is real physics.} Between the onset and the
point where the deconfined branch turns around, $\mathrm dP/\mathrm dn_B<0$:
the mechanically unstable side. A raw branch may carry it; a construction
(Maxwell, Gibbs, or the $\eta$-mixed phase) removes it before any table reaches
a structure solver. Sound speeds must be taken one-sided on each branch, never
across the transition.

\section{Modes}
\label{sec:modes}

The closure rows of the specification map onto the repository's four modes:

\begin{table}[h]
\centering
\begin{tabular}{lllp{6.2cm}}
\toprule
row & mode & variables & conditions \\
\midrule
R1 & \texttt{beta\_eq\_neutrinoless} & $(n_B,T)$ &
 $n_C=n_e+n_\mu$; $\mu_e=-\mu_C$; $\mu_S=0$ \\
R3 & \texttt{beta\_eq\_neutrino\_trapped} & $(n_B,Y_{L_e},T)$ &
 as R1 plus $n_{L_e}=Y_{L_e}n_B$; $\nu$ with $g=1$ \\
R2 & \texttt{fixed\_YC}, \texttt{leptons=True} & $(n_B,Y_C,T)$ &
 $n_C=Y_Cn_B$ and $n_C=n_e+n_\mu$; \textbf{no} weak equilibrium \\
R4 & \texttt{fixed\_YC\_YS}, \texttt{leptons=False} & $(n_B,Y_C,Y_S,T)$ &
 $n_C=Y_Cn_B$, $n_S=Y_Sn_B$; no leptons \\
R5 & \texttt{fixed\_YC\_YS} at $Y_C=\tfrac12,Y_S=0$ & $(n_B,T)$ &
 symmetric matter \\
\bottomrule
\end{tabular}
\end{table}

R2 is the one to state twice: it imposes total electric neutrality
\emph{without} imposing beta equilibrium, which is what merger and
core-collapse matter is on a dynamical timescale. \textbf{Weak equilibrium is a
per-row closure, never an identity built into $\Omega$}; hardwiring it would
make such matter unrepresentable. Fractions are per baryon, $Y_C=n_C/n_B$ and
$Y_S=n_S/n_B$, which differ by a factor three from the per-quark
$n_s/(3n_B)$ often plotted.

Wherever a temperature is accepted, entropy per baryon is accepted in its
place, through an outer one-dimensional solve for $T$.

\section{Parameters}
\label{sec:parameters}

\subsection{Fixed by vacuum data}

$f_\pi=93$, $m_\pi=138$, $f_K=113$, $m_K=496$, $m_{u,d}=5$, $m_s=95$,
$m_\zeta=980$, $m_\varphi=1600$ (lattice scalar glueball), $m_\omega=783$ MeV.
Everything below follows in closed form:
\begin{align}
\sigma_0 &= f_\pi, &
\zeta_0 &= \sqrt2 f_K - \frac{f_\pi}{\sqrt2}, &
\varphi_0 &= \frac{4\sqrt{B_g}}{m_\varphi}, \\
\epsilon_\sigma &= f_\pi m_\pi^2, &
\epsilon_\zeta &= \sqrt2 f_K m_K^2 - \frac{f_\pi}{\sqrt2}m_\pi^2, & & \\
\lambda &= \frac{m_\sigma^2-m_\pi^2}{2f_\pi^2}, &
v^2 &= f_\pi^2 - \frac{m_\pi^2}{\lambda}, & & \\
\lambda_\zeta &= \frac{m_\zeta^2-\epsilon_\zeta/\zeta_0}{2\zeta_0^2}, &
v_\zeta^2 &= \zeta_0^2 - \frac{\epsilon_\zeta}{\lambda_\zeta\zeta_0}, & &
\end{align}
and $C_0$ from $V(\sigma_0,\zeta_0)=0$. Numerically at $m_\sigma=550$ MeV:
$\zeta_0=94.045$ MeV, $\lambda=16.387$, $v=86.527$ MeV,
$\lambda_\zeta=31.414$, $v_\zeta^2=-4039.3$ MeV$^2$,
$C_0=2.4352\times10^{9}$ MeV$^4$, $\varphi_0=56.25$ MeV.

\textbf{$v_\zeta^2$ is negative} at the baseline $m_\zeta$: the strange quartic
is convex, explicit breaking dominating, so the strange sector does not break
chirally on its own in this truncation. The sign flips between $m_\zeta=1100$
and $1150$ MeV. Never assume it is positive and never write $v_\zeta$ as a
square root.

\subsection{The effective bag constant}

Not an input. At the physical vacuum the field energy is zero by construction;
at the perturbative point it is
\begin{equation}
B_{\rm eff} = \underbrace{U(0)-U(\varphi_0)}_{B_g}
 + \underbrace{V(0,0)-V(f_\pi,\zeta_0)}_{B_\chi}.
\end{equation}
At $B_g^{1/4}=150$, $m_\sigma=550$, $m_\zeta=980$ MeV this gives
$B_\chi^{1/4}=229.9$ MeV and
$B_{\rm eff}=(239.7\ \mathrm{MeV})^4=429.4$ MeV\,fm$^{-3}$.
\textbf{The chiral sector supplies the larger part}: quoting $B_g$ alone as
"the bag constant" of this model is wrong by a factor of six in energy density.

\subsection{Free parameters}

$B_g^{1/4}$ (120--250 MeV), $g_q$ (3--6), $g_s$ (3--8), $m_\sigma$
(450--700 MeV), $\bar g_\omega$ (0--12), $n_c$ (0.3--3 fm$^{-3}$), and at L3
$G_D$ and $\Lambda$ (550--800 MeV, nearly degenerate with $G_D$). The discrete
choices $p=1$ and $q\in\{0,1\}$ are declared per run, not sampled.
$m_\omega$ is a normalisation convention, only $g_\omega/m_\omega$ entering;
$m_\varphi$ cancels from the bulk equation of state at fixed $B_g$, pricing
gradients and the glueball spectrum instead.

$g_q=3.0$ is pinned by the specification's \S10 table, which quotes
$\Mstar_{u,d}=826$ MeV at $\phib=0.90$ and $1531$ MeV at $\phib=0.95$ in the
confined branch; both invert to $g_q=3.00$. $g_s$, $\bar g_\omega$ and $n_c$
are calibration knobs and the shipped values are mid-prior rather than
measurements; the code says so. $G_D=5\times10^{-6}$ MeV$^{-2}$ is calibrated
here so that the gap falls inside the specification's 20--150 MeV window at
$\mu_q\simeq450$ MeV.

\section{What is not in the model}

The dilaton gradient and finite-size terms (this is a bulk, homogeneous mean
field). The de Carvalho contact coupling $h(\varphi)$ as a replacement for
$G_D$: their construction is sound --- integrating out the confining field's
fluctuation at Gaussian order gives a contact interaction, which is why
$\mathcal L_{\rm pair}$ is a legitimate leading term rather than an ad hoc
addition, and it gives "zero range" the physical meaning $1/M_\chi$ --- but
what it predicts is a coupling negligible where the quarks are light and
enormous where they are heavy, which removes pairing from the deconfined phase
where a star's core lives and would condense a diquark in the confining vacuum.
It carries $q=4$ at $p=1$, violating $q\le p$. What is taken from it is the
argument, not the number.

\section{References}

\begin{itemize}
\item Friedberg, Lee, \emph{Phys. Rev. D} \textbf{15} (1977) 1694;
 \textbf{16} (1977) 1096; \textbf{18} (1978) 2623.
\item Nielsen, Patk\'os, \emph{Nucl. Phys. B} \textbf{195} (1982) 137.
\item Mack, \emph{Nucl. Phys. B} \textbf{235} (1984) 197.
\item Pirner, \emph{Prog. Part. Nucl. Phys.} \textbf{29} (1992) 33.
\item Birse, \emph{Prog. Part. Nucl. Phys.} \textbf{25} (1990) 1.
\item Drago, Fiolhais, Tambini, arXiv:hep-ph/9503462.
\item Ghosh, Phatak, \emph{J. Phys. G} \textbf{18} (1992) 755;
 \emph{Phys. Rev. C} \textbf{52} (1995) 2195, arXiv:nucl-th/9509017.
\item Alberico, Drago, Ratti, arXiv:hep-ph/0110091.
\item Maieron, Baldo, Burgio, Schulze, arXiv:nucl-th/0404089.
\item de Carvalho, Malheiro et al., \emph{Nucl. Phys. B Proc. Suppl.}
 \textbf{199} (2010) 308.
\item Buballa, \emph{Phys. Rept.} \textbf{407} (2005) 205, arXiv:hep-ph/0402234.
\item Alford, Schmitt, Rajagopal, Sch\"afer, \emph{Rev. Mod. Phys.}
 \textbf{80} (2008) 1455, arXiv:0709.4635.
\item Steiner, Reddy, Prakash, \emph{Phys. Rev. D} \textbf{66} (2002) 094007,
 arXiv:hep-ph/0205201.
\item R\"uster, Werth, Buballa, Shovkovy, Rischke, \emph{Phys. Rev. D}
 \textbf{72} (2005) 034004, arXiv:hep-ph/0503184.
\item Typel, Oertel, Kl\"ahn et al., CompOSE manual, arXiv:2203.03209.
\end{itemize}

\end{document}
